\scanpage{41}

compatibles avec l'action de $\Pi_\infty$ et de $\unsure{\mathfrak S_3|\mathscr S_3}\deleted{\unsure{\Gamma}}\times\underline{\{1,\tau\}}$ ($\tau\inserted{\in\Gamma}$ étant la conjugaison complexe) il suffit pour déterminer ces $18$ flèches, de déterminer deux d'entre elles, \deleted{telles que} \inserted{p. ex}

\[
E(\underset{\substack{\Vert\\\alpha_1}}{\alpha})\longrightarrow
\underset{\substack{\Vert\\E_{1/2}}}{E'_\infty}\longrightarrow E_0
\]

[NB $E(\alpha)\mathrel{\unsure{-|\longrightarrow}}E_0$ est le composé] et $=$ de connaître \inserted{les} images dans $E_0,E_1$ de l'élément marqué $\unsure{\sigma_\infty|\alpha_\infty}$ de $E(\sigma)$, [d'où élément marqué $\deleted{\beta}\alpha'_\infty(i)\in E'_i=\deleted{\unsure{/|,}}E(\alpha'(i))$, et $\deleted{E}\alpha_\infty(i)\in E_i=E(\alpha(i))$, $\forall i\in\{0,1,\infty\}$].

On trouve

\[
\begin{aligned}
\alpha'_n(\infty)&=\left(\sqrt[n]{\tfrac12}\exp\tfrac{-\pi i}{n},\sqrt[n]{\tfrac12}\exp\tfrac{+\pi i}{n},1\right)\\
&=(\deleted{\unsure{\varphi}}\theta_n^{-1}\bar\zeta_n,\theta_n^{-1}\zeta_n,1)\\
\deleted{\alpha'_\infty(\infty)}&\deleted{=}\\[-.6ex]
\alpha'_\infty(\infty)&=(\theta_\infty\wedge\bar\zeta_{\unsure{n|\infty}},\theta_\infty\wedge\zeta_\infty,1)
\end{aligned}
\]

% archive sheet number: 40
